\documentclass{amsart}
% For dealing with references we use the comment environment
\usepackage{verbatim}
\newenvironment{reference}{\comment}{\endcomment}
%\newenvironment{reference}{}{}
\newenvironment{slogan}{\comment}{\endcomment}
\newenvironment{history}{\comment}{\endcomment}

% For commutative diagrams we use Xy-pic
\usepackage[all]{xy}

% We use 2cell for 2-commutative diagrams.
\xyoption{2cell}
\UseAllTwocells

% We use multicol for the list of chapters between chapters
\usepackage{multicol}

% This is generally recommended for better output
\usepackage{lmodern}
\usepackage[T1]{fontenc}

% For cross-file-references

% Package for hypertext links:
\usepackage{hyperref}

% For any local file, say "hello.tex" you want to link to please

% Theorem environments.
%
\theoremstyle{plain}
\newtheorem{theorem}[subsection]{Theorem}
\newtheorem{proposition}[subsection]{Proposition}
\newtheorem{lemma}[subsection]{Lemma}

\theoremstyle{definition}
\newtheorem{definition}[subsection]{Definition}
\newtheorem{example}[subsection]{Example}
\newtheorem{exercise}[subsection]{Exercise}
\newtheorem{situation}[subsection]{Situation}

\theoremstyle{remark}
\newtheorem{remark}[subsection]{Remark}
\newtheorem{remarks}[subsection]{Remarks}

\numberwithin{equation}{subsection}

% Macros
%
\def\lim{\mathop{\mathrm{lim}}\nolimits}
\def\colim{\mathop{\mathrm{colim}}\nolimits}
\def\Spec{\mathop{\mathrm{Spec}}}
\def\Hom{\mathop{\mathrm{Hom}}\nolimits}
\def\Ext{\mathop{\mathrm{Ext}}\nolimits}
\def\SheafHom{\mathop{\mathcal{H}\!\mathit{om}}\nolimits}
\def\SheafExt{\mathop{\mathcal{E}\!\mathit{xt}}\nolimits}
\def\Sch{\mathit{Sch}}
\def\Mor{\mathop{\mathrm{Mor}}\nolimits}
\def\Ob{\mathop{\mathrm{Ob}}\nolimits}
\def\Sh{\mathop{\mathit{Sh}}\nolimits}
\def\NL{\mathop{N\!L}\nolimits}
\def\CH{\mathop{\mathrm{CH}}\nolimits}
\def\proetale{{pro\text{-}\acute{e}tale}}
\def\etale{{\acute{e}tale}}
\def\QCoh{\mathit{QCoh}}
\def\Ker{\mathop{\mathrm{Ker}}}
\def\Im{\mathop{\mathrm{Im}}}
\def\Coker{\mathop{\mathrm{Coker}}}
\def\Coim{\mathop{\mathrm{Coim}}}

% Boxtimes
%
\DeclareMathSymbol{\boxtimes}{\mathbin}{AMSa}{"02}

%
% Macros for moduli stacks/spaces
%
\def\QCohstack{\mathcal{QC}\!\mathit{oh}}
\def\Cohstack{\mathcal{C}\!\mathit{oh}}
\def\Spacesstack{\mathcal{S}\!\mathit{paces}}
\def\Quotfunctor{\mathrm{Quot}}
\def\Hilbfunctor{\mathrm{Hilb}}
\def\Curvesstack{\mathcal{C}\!\mathit{urves}}
\def\Polarizedstack{\mathcal{P}\!\mathit{olarized}}
\def\Complexesstack{\mathcal{C}\!\mathit{omplexes}}
% \Pic is the operator that assigns to X its picard group, usage \Pic(X)
% \Picardstack_{X/B} denotes the Picard stack of X over B
% \Picardfunctor_{X/B} denotes the Picard functor of X over B
\def\Pic{\mathop{\mathrm{Pic}}\nolimits}
\def\Picardstack{\mathcal{P}\!\mathit{ic}}
\def\Picardfunctor{\mathrm{Pic}}
\def\Deformationcategory{\mathcal{D}\!\mathit{ef}}

\setlength{\emergencystretch}{3em}
\begin{document}
\title[Stacks Project, Tag 09II]{423 --- Elimination of ramification for essentially finite-type valuation-ring extensions}
\maketitle
\noindent Copyright (C) 2005--2025 Johan de Jong. Stacks Project authors. Source: \href{https://stacks.math.columbia.edu/tag/09II}{Tag 09II}.

\noindent Editorial difficulty note (not author text). Difficulty H3: The author perfects the residue field through a filtered union of weakly unramified extensions and combines Epp solutions for all resulting branches. It then descends the solution to a finite field-extension stage by stabilizing spectra, uniformizers and separable residue-field generators..

\noindent Editorial provenance note (not author text). History: excerpt from revision \texttt{a0444\allowbreak 6e57e\allowbreak c1fbc\allowbreak 252a8\allowbreak 71afc\allowbreak ec775\allowbreak 2fb28\allowbreak 07b14}, more-algebra.tex, lines 35960--36069. Reference presentation was adapted; original-author context and core support blocks were retained or added. The mathematical wording is unchanged. Licensed under GNU FDL 1.2 or later; no invariant sections or cover texts. See ../licenses/COPYING. Context blocks reproduce author definitions and supporting results; they are not additional counted problems.

\begin{definition}
\label{definition-solution}
Let $A \to B$ be an extension of discrete valuation rings with fraction
fields $K \subset L$.
\begin{enumerate}
\item We say a finite field extension $K_1/K$ is a
{\it weak solution for $A \subset B$} if all the extensions
$(A_1)_{\mathfrak m_i} \subset (B_1)_{\mathfrak m_{ij}}$ of
Remark \href{https://stacks.math.columbia.edu/tag/09EM}{remark construction (Tag 09EM)} are weakly unramified.
\item We say a finite field extension $K_1/K$ is a
{\it solution for $A \subset B$} if each extension
$(A_1)_{\mathfrak m_i} \subset (B_1)_{\mathfrak m_{ij}}$ of
Remark \href{https://stacks.math.columbia.edu/tag/09EM}{remark construction (Tag 09EM)} is formally smooth in
the $\mathfrak m_{ij}$-adic topology.
\end{enumerate}
We say a solution $K_1/K$ is a {\it separable solution}
if $K_1/K$ is separable.
\end{definition}

\begin{proposition}
\label{proposition-epp-essentially-finite-type}
\begin{reference}
See \href{https://stacks.math.columbia.edu/bibliography}{Bibliography: alterations (Lemma 2.13)} for a special case.
\end{reference}
Let $A \to B$ be an extension of discrete valuation rings with fraction
fields $K \subset L$. If $B$ is essentially of finite type over $A$, then
there exists a finite extension $K_1/K$ which is a solution for
$A \to B$ as defined in
Definition \ref{definition-solution}.
\end{proposition}

\begin{proof}
Observe that a weak solution is a solution if the residue field of $A$
is perfect, see Lemma \href{https://stacks.math.columbia.edu/tag/09E7}{lemma extension dvrs formally smooth (Tag 09E7)}.
Thus the proposition follows immediately from Theorem \href{https://stacks.math.columbia.edu/tag/09F9}{theorem epp (Tag 09F9)}
if the residue characteristic of $A$ is $0$ (and in fact we do not need
the assumption that $A \to B$ is essentially of finite type).
If the residue characteristic of $A$ is $p > 0$ we will also deduce it from
Epp's theorem.

\medskip\noindent
Let $x_i \in A$, $i \in I$ be a set of elements mapping to a $p$-base of
the residue field $\kappa$ of $A$. Set
$$
A' = \bigcup\nolimits_{n \geq 1} A[t_{i, n}]/(t_{i, n}^{p^n} - x_i)
$$
where the transition maps send $t_{i, n + 1}$ to $t_{i, n}^p$. Observe
that $A'$ is a filtered colimit of weakly unramified finite extensions
of discrete valuation rings over $A$. Thus $A'$ is a discrete valuation
ring and $A \to A'$ is weakly unramified. By construction the residue field
$\kappa' = A'/\mathfrak m_A A'$ is the perfection of $\kappa$.

\medskip\noindent
Let $K'$ be the fraction field of $A'$.
We may apply Lemma \href{https://stacks.math.columbia.edu/tag/09IH}{lemma big extension is ok (Tag 09IH)}
to the extension $K'/K$. Thus $B'$ is a finite product of
Dedekind domains. Let $\mathfrak m_1, \ldots, \mathfrak m_n$ be the
maximal ideals of $B'$. Using Epp's theorem (Theorem \href{https://stacks.math.columbia.edu/tag/09F9}{theorem epp (Tag 09F9)})
we find a weak solution $K'_i/K'$ for each of the
extensions $A' \subset B'_{\mathfrak m_i}$. Since the residue field
of $A'$ is perfect, these are actually solutions. Let $K'_1/K'$
be a finite extension which contains each $K'_i$. Then $K'_1/K'$
is still a solution for each $A' \subset B'_{\mathfrak m_i}$ by
Lemma \href{https://stacks.math.columbia.edu/tag/0GLR}{lemma solution goes up (Tag 0GLR)}.

\medskip\noindent
Let $A'_1$ be the integral closure of $A$ in $K'_1$. Note that
$A'_1$ is a Dedekind domain by the discussion in
Remark \href{https://stacks.math.columbia.edu/tag/09EM}{remark construction (Tag 09EM)} applied to $K' \subset K'_1$.
Thus Lemma \href{https://stacks.math.columbia.edu/tag/09IH}{lemma big extension is ok (Tag 09IH)} applies to $K'_1/K$.
Therefore the integral closure $B'_1$ of $B$ in
$L'_1 = (L \otimes_K K'_1)_{red}$ is a Dedekind domain and because
$K'_1/K'$ is a solution for each $A' \subset B'_{\mathfrak m_i}$
we see that $(A'_1)_{A'_1 \cap \mathfrak m} \to (B'_1)_{\mathfrak m}$
is formally smooth in the $\mathfrak m$-adic topology
for each maximal ideal $\mathfrak m \subset B'_1$.

\medskip\noindent
By construction, the field $K'_1$ is a filtered colimit of finite
extensions of $K$. Say $K'_1 = \colim_{i \in I} K_i$. For each $i$ let
$A_i$, resp.\ $B_i$ be the integral closure of
$A$, resp.\ $B$ in $K_i$, resp.\ $L_i = (L \otimes_K K_i)_{red}$.
Then it is clear that
$$
A'_1 =  \colim A_i\quad\text{and}\quad B'_1 = \colim B_i
$$
Since the ring maps $A_i \to A'_1$ and $B_i \to B'_1$ are injective
integral ring maps and since $A'_1$ and $B'_1$ have finite spectra,
we see that for all $i$ large enough the ring maps
$A_i \to A'_1$ and $B_i \to B'_1$ are bijective on spectra.
Once this is true, for all $i$ large enough the maps
$A_i \to A'_1$ and $B_i \to B'_1$ will be weakly unramified
(once the uniformizer is in the image). It follows from multiplicativity
of ramification indices that $A_i \to B_i$ induces weakly unramified maps
on all localizations at maximal ideals of $B_i$ for such $i$.
Increasing $i$ a bit more we see that
$$
B_i \otimes_{A_i} A'_1 \longrightarrow B'_1
$$
induces surjective maps on residue fields (because the residue fields
of $B'_1$ are finitely generated over those of $A'_1$ by
Lemma \href{https://stacks.math.columbia.edu/tag/09IH}{lemma big extension is ok (Tag 09IH)}). Picture of residue
fields at maximal ideals lying under a chosen maximal ideal
of $B'_1$:
$$
\xymatrix{
\kappa_{B_i} \ar[r] &
\kappa_{B_{i'}} \ar[r] &
 & \ldots &
\kappa_{B'_1} \\
\kappa_{A_i} \ar[r] \ar[u] &
\kappa_{A_{i'}} \ar[r] \ar[u] &
 & \ldots &
\kappa_{A'_1} \ar[u]
}
$$
Thus $\kappa_{B_i}$ is a finitely generated extension of
$\kappa_{A_i}$ such that the compositum of $\kappa_{B_i}$
and $\kappa_{A'_1}$ in $\kappa_{B'_1}$ is separable over
$\kappa_{A'_1}$. Then that happens already at a finite stage:
for example, say $\kappa_{B'_1}$ is finite separable over
$\kappa_{A'_1}(x_1, \ldots, x_n)$, then just increase $i$
such that $x_1, \ldots, x_n$ are in $\kappa_{B_i}$ and such that
all generators satisfy separable polynomial equations over
$\kappa_{A_i}(x_1, \ldots, x_n)$. This means that
$A_i \to (B_i)_\mathfrak m$ is formally smooth in the $\mathfrak m$-adic
topology for all maximal ideals $\mathfrak m$ of
$B_i$ and the proof is complete.
\end{proof}
\end{document}
